\centerline{\bf \Huge Сравнения по модулю I}

\textbf{Определение 1.} Числа $a$ и $b$ \textit{сравнимы по модулю} $m$, если $a-b$ делится на $m$.

\textbf{Определение 2.} Числа $a$ и $b$ \textit{сравнимы по модулю} $m$, если $a$ и $b$ имеют одинаковые остатки при делении на $m$.

\textbf{Определение.} Классом вычетов по данному модулю $m$ называется множество всех целых чисел сравнимых с некоторым данным целым числом $a$ по модулю $m$. Такой класс обозначается $\overline{\mathrm{a}}$.

\centerline{\it Свойства сравнений.}

1. $a \equiv a\ (mod\ m)$.

2. Если $a \equiv b\ (mod\ m)$, то $b \equiv a\ (mod\ m)$.

3. Если $a \equiv b\ (mod\ m)$ и $b \equiv c\ (mod\ m)$, то $a \equiv c\ (mod\ m)$.

4. Если $a \equiv b\ (mod\ m)$ и $c \equiv d\ (mod\ m)$, то $a \pm c \equiv b \pm d\ (mod\ m)$.

5. Если $a \equiv b\ (mod\ m)$, то $a c \equiv b c\ (mod\ m)$.

6. Если $a \equiv b\ (mod\ m)$ и $c \equiv d\ (mod\ m)$, то $a c \equiv b d\ (mod\ m)$.

7. Если $a \equiv b\ (mod\ m)$, то $a^n \equiv b^n\ (mod\ m)$.

8. Если $k a \equiv k b\ (mod\ m)$ и НОД $(k, m)=1$, то $a \equiv b\ (mod\ m)$.

\problem{1}
Найдите остаток от деления

а) $7778 \cdot 7779 \cdot 7780 \cdot 7781 \cdot 7782 \cdot 7783$ на 7 ;

б) $2013 \cdot 2014 \cdot 2015 \cdot 2016$ на 2017 ;

в) $1 \cdot 3 \cdot 5 \ldots 101+2 \cdot 4 \cdot 6 \ldots 102$ на 103 .


\problem{2}
Докажите, что при любом натуральном $n$ число $16^{n+2}+23^{n+1}+37^n$ делится на 7.

\begin{flushright}
        \scriptsize{Без МТФ}
\end{flushright}

\problem{3}
Числа $m$ и $n$ нечётны. Докажите, что $1^n+2^n+3^n+\ldots+(m-1)^n \ \vdots \ m$.

\newpage

\hspace{3mm}

\problem{4}
\textbf{Малая теорема Ферма.}

a) Дано натуральное число $n$. При каких натуральных $a$ среди чисел $0 \cdot a, 1 \cdot a, 2 \cdot a, \ldots,(n-1) \cdot a$ встречаются все остатки при делении на $n ?$

б) Докажите, что если $(a, p) = 1$, то $a^{p-1} \equiv 1 \ (mod \ p)$, где $p$ простое число

\problem{5}
\textbf{Существование <<обратного элемента>>.} Выведите из задачи $4a)$, что у сравнения $ax = 1 \ (mod \ p)$ существует не более одного решения, где $p$ простое число.

\problem{6}
\textbf{Теорема Вильсона.} Выведите из задачи $5$, что $(n-1)! \equiv 1 \ (mod \ n)$ тогда и только тогда, когда $n$ --- простое число.